\subsection{An adjunction obstruction to a separable-cover construction}
\paragraph{A2Z (Proved exclusion of the stated construction).}
Let $k$ be algebraically closed of characteristic $p>0$, and let
$\pi:X\to B$ be a finite, flat, generically separable morphism of
degree $p$ between smooth projective integral varieties of dimension
$n\geq2$. Let $C\subset B$ be a smooth ample integral divisor, and
assume that its scheme-theoretic inverse image
\[
 Y=\pi^*C
\]
is smooth and integral. Then
\[
 \operatorname{Pic}(C)[p](k)\longrightarrow
 \operatorname{Pic}(Y)[p](k)
\]
is injective. In particular this construction cannot make a
nontrivial $M|_C$, for $M\in\operatorname{Pic}(B)[p](k)$, disappear
on $Y$. This is a restriction on the proposed construction; it does
not prove the general point-injectivity assertion A2.

Suppose that $N$ is a nontrivial order-$p$ line bundle on $C$ and
$\rho^*N$ is trivial, where $\rho=\pi|_Y$. Choose $N^p\simeq
\mathcal O_C$, and let $\tau:T\to C$ be the associated
$\mu_p$-torsor. It is integral. Indeed, if its generic equation
$z^p=u$ split in $k(C)$, the rational root would give a rational
trivialization of $N$ with zero divisor, hence a global
trivialization because $C$ is normal. Its degree is $p$.
A trivialization of $\rho^*N$ can be made compatible with the
chosen $p$th-power trivialization by rescaling by a constant:
$H^0(Y,\mathcal O_Y^*)=k^*$ and $k$ is algebraically closed.
It therefore gives a factorization
\[
 Y\xrightarrow{\nu}T\xrightarrow{\tau}C.
\]
Since both $\rho$ and $\tau$ have degree $p$, $\nu$ is finite
birational. Thus $Y$ is the normalization of $T$. In particular
$\rho$ is purely inseparable.

The torsor $T$ is Gorenstein and
\[
 \omega_T\simeq\tau^*\omega_C.                         \tag{A2Z.1}
\]
For an explicit check, realize it inside the total space of the
line bundle $N$ by its equation $z^p=u$. Adjunction gives
$\omega_T\simeq\tau^*(\omega_C\otimes N^{p-1})$; the
nowhere-zero tautological root trivializes $\tau^*N$, giving
(A2Z.1). This argument applies when $T$ is nonnormal.

Write $\mathfrak c$ for the conductor, viewed as an ideal of
$\mathcal O_Y$. Finite duality identifies
\[
 \omega_Y\simeq\mathfrak c\otimes\nu^*\omega_T.
\]
Here the conductor identification is literal: on an affine
birational extension $A\subset A'$ inside a common fraction field,
\[
 \operatorname{Hom}_A(A',A)
   =\{q:qA'\subset A\}
   =\operatorname{Ann}_A(A'/A).
\]
Both canonical sheaves are invertible, so $\mathfrak c$ is an
invertible nonzero ideal on $Y$. Consequently it is
$\mathcal O_Y(-D)$ for an effective Cartier divisor $D$, and
\[
 \omega_Y\simeq\rho^*\omega_C\otimes\mathcal O_Y(-D).
                                                               \tag{A2Z.2}
\]
The finite-duality input is the $R\operatorname{Hom}$ description
of the dualizing complex under a finite map
\cite[Lemma 47.15.8, with the finite-module description in Section 49.5]{DGJ_StacksFiniteDuality};
taking the common lowest cohomological degree gives the displayed
module identity. The divisor adjunction used here also follows
from \cite[Lemma 48.14.1, Tag 0B4B]{DGJ_StacksAdjunction}.

On the other hand, the determinant of the generically invertible
map $d\pi:\pi^*\Omega_B^1\to\Omega_X^1$ is a nonzero regular
section of $\omega_X\otimes(\pi^*\omega_B)^{-1}$. Let $R$ be
its effective divisor. Hence
\[
 \omega_X\simeq\pi^*\omega_B\otimes\mathcal O_X(R).
                                                               \tag{A2Z.3}
\]
The divisor $R$ contains $Y$. At its generic point the normal
map in the conormal sequences for $Y\subset X$ and $C\subset B$
is an isomorphism, because $\pi^*C=Y$ scheme-theoretically,
whereas the tangential determinant vanishes since $\rho$ is
purely inseparable of degree $p$. Thus the full determinant
vanishes at that generic point.
Applying smooth-divisor adjunction to (A2Z.3) gives
\[
 \omega_Y\simeq\rho^*\omega_C\otimes
                  \mathcal O_X(R)|_Y.                         \tag{A2Z.4}
\]
Comparison with (A2Z.2) says
$\mathcal O_X(R)|_Y\otimes\mathcal O_Y(D)\simeq\mathcal O_Y$.
Intersecting with $\mathcal O_X(Y)|_Y^{\,n-2}$ gives
\[
 0=R\cdot Y^{n-1}+D\cdot (Y|_Y)^{n-2}>0,
\]
a contradiction: $Y$ is ample because it is a finite pullback
of $C$, $R$ is a nonzero effective divisor (indeed $R\geq Y$),
and $D$ is effective. The first summand is strictly positive and
the second is nonnegative. This completes the proof.

\emph{Scope of the obstruction.}
A local separable equation such as
$z^p-s^{p-1}z=b$ can have a smooth total space and purely
inseparable residue extension at $s=0$. The proposition explains
why these local equations cannot be glued into the stated
projective construction with a smooth ample full inverse image
that kills a nontrivial $p$-torsion class. Allowing the
$\mu_p$-torsor on $C$ to be nonnormal does not evade the argument:
its conductor contributes the additional nonnegative term above.
The degree-$p$ hypothesis and the equality $\pi^*C=Y$ are retained;
no assertion about arbitrary higher-degree covers or a single
component of a reducible inverse image is made.
